\question
После долгих размышлений у печки и научного спора с Дядей Фёдором рациональный кот Матроскин решил подойти к вопросу бутербродоедения с истинно научной строгостью. Он скрупулёзно записывал сведения о бутербродах и их вкусовых качествах. Результатом наблюдений стало то,
что сэндвич считается идеальным при следующем раскладе: (колбаса сверху И хлеб свежий И (масло XOR икра)) ИЛИ (колбаса снизу И (свежий хлеб ИЛИ масло) И НЕ хлеб чёрствый). В остальных случаях бутерброд получался "так себе". 

\textit{Замечание:} считайте, что «колбаса снизу» и «хлеб чёрствый» --- это строгие логические отрицания событий «колбаса сверху» и «хлеб свежий».

Матроскин недавно увидел статью про карты Вейча как способ минимизации функций, но, к сожалению, ничего не понял, поэтому решил попросить вашей помощи.

Превратите эти гастрономические наблюдения в логическую формулу, да не простую, а компактную (найти минимальную ДНФ), чтобы Матроскин мог быстро определять оптимальный способ бутербродоукладывания.

---------------

Автор -- Никита Глейм